\item \textbf{Problema de repaso.} Considere un objeto rígido, homogéneo, cilíndrico, lineal, superficial o esférico. ¿Cuál es el aumento porcentual en el momento de inercia del objeto cuando se calienta de 0 a 100 °C si está compuesto de (a) cobre o (b) aluminio? Suponga que los coeficientes de expansión lineal promedio que se muestran en la tabla 6.1 no varían entre 0 y 100 °C. (c) ¿Por qué las respuestas a los incisos (a) y (b) son las mismas para todas las formas?
